\documentclass[11pt,a4paper]{article}

\usepackage[T1]{fontenc}
\usepackage[utf8]{inputenc}
\usepackage{lmodern}
\usepackage[margin=1in]{geometry}
\usepackage{microtype}
\usepackage{amsmath}
\usepackage{amssymb}
\usepackage{booktabs}
\usepackage{graphicx}
\usepackage{xcolor}
\usepackage{caption}
\usepackage{natbib}
\usepackage[hidelinks]{hyperref}

\captionsetup{font=small,labelfont=bf}

\title{%
  \textbf{OMAT-SE}\\[0.35em]
  \large A Faceted Search Engine for Materials Property Data
}
\author{Darshan Jaidka}
\date{\today}

\begin{document}

\maketitle

\begin{abstract}
Materials property data is scattered across journal tables and published as
unstructured text, which makes comparison across compositions and applications
slow. This report describes OMAT-SE, a search system built around a curated
corpus of measurement records. Each record links a paper to an elemental
composition, a measured property with its unit, a journal, a table reference,
and an application group. The system combines BM25 retrieval, faceted filtering,
and aggregation-backed visualisations, and exposes the same domain layer through
a REST API and a Model Context Protocol server. A per-paper chat interface
answers questions from the paper's own measurement rows and cites the records it
uses. This document defines the data contract, the architecture, the ranking and
filtering design, and the evaluation plan.
\end{abstract}

\section{Introduction}

Searching materials literature usually means recalling a composition, guessing a
property name, and scanning tables by hand. The underlying data is relational:
papers report compositions; compositions have measured properties; properties are
associated with applications. A useful search engine therefore has to support
three different entry points at once: an exact material label, an elemental
constraint such as ``contains Al and Cr'', and a numeric range such as ``Young's
modulus above 400~GPa''.

Structured efforts such as the Materials Project demonstrate the value of
queryable property data \citep{jain2013}, yet a large share of published
measurements remains locked in journal tables. OMAT-SE treats the measurement as
the primary entity. A paper is a grouping and a provenance record, not the unit
of retrieval. This choice keeps facet counts honest and makes numeric filters
cheap. The contributions described here are
(i) a canonical data contract for the source CSV, (ii) an OpenSearch index and
query design for faceted materials search, (iii) a per-paper chat interface that
does not require a retrieval-augmented generation pipeline, and (iv) a
reproducible evaluation plan.

\section{Data}

The source is a CSV with eleven columns. Composition, property, element, and
application cells are Python literals, so parsing uses \texttt{ast.literal\_eval}
rather than regular expressions. Table~\ref{tab:columns} summarises the mapping
from source columns to canonical fields.

\begin{table}[htbp]
  \centering
  \small
  \begin{tabular}{@{}llp{0.42\linewidth}@{}}
    \toprule
    Source column & Canonical field & Treatment \\
    \midrule
    Article PII & \texttt{article\_pii} & Publisher Item Identifier; keyword \\
    DOI & \texttt{doi} & Lowercased, prefix stripped; keyword \\
    Table No. & \texttt{table\_nos} & Integer list; provenance in the UI \\
    ID & \texttt{legacy\_id} & Source label; checked for cross-paper reuse \\
    Proxy\_ID & \texttt{proxy\_id} & Primary record key \\
    Composition & \texttt{composition} & Label plus authoritative elemental amounts \\
    Property & \texttt{property} & Name, value, unit; one document per property \\
    Journal\_Name & \texttt{journal} & Underscores expanded; slug facet \\
    Elements & \texttt{elements} & Validated against the composition amounts \\
    APL\_Group & \texttt{applications} & Application taxonomy and slug \\
    Date & \texttt{year} & Integer range filter and histogram \\
    \bottomrule
  \end{tabular}
  \caption{Source columns and their canonical representation. \texttt{Paper\_Link}
  is derived from the DOI and is not present in the source.}
  \label{tab:columns}
\end{table}

A sample record has composition label \texttt{TiAl70CrN} with measured amounts
Al~95.89, Cr~1.37, N~1.37, and Ti~1.37 in mol\%, and a Young's modulus of
467~GPa. The label is not a stoichiometric formula: the amounts are the
authoritative representation. This distinction matters because a formula parser
would silently produce a composition that contradicts the measurement.

\subsection{Cleaning and validation}

Cells are parsed into typed values and validated before indexing. Rejected rows
are written to a rejects directory with the original line and a machine-readable
reason. The rules are deliberately conservative: elemental amounts must sum to
$100 \pm 0.5$; the element set must match the composition list; units must exist
in a curated unit table; years must fall between 1900 and 2027; and every DOI
must match the standard prefix pattern. DOIs also provide a stable path to
publisher metadata and landing pages \citep{hendricks2020}. Rows with an empty
application group are labelled \texttt{unclassified} rather than dropped.

Property names and journals are normalised through curated YAML taxonomies.
These files are reviewed as data, not inferred in code, because a wrong mapping
is a silent search failure. Raw labels are always preserved next to normalised
ones so that users can verify a record against the source table.

\section{Architecture}

The system is split into five deployable units. A Next.js application renders
search and paper pages on the server and is deployed on Vercel. A FastAPI service
validates requests, applies caching, and shapes responses. A shared Python package
contains the domain models, cleaning rules, and OpenSearch query builders. An
ingestion CLI performs validation and bulk indexing. A small Model Context
Protocol server exposes the same search tools to agents.

API requests pass through CloudFront and an application load balancer. Redis
caches first-page results, facet counts, and autocomplete suggestions. OpenSearch
is never exposed publicly. PostgreSQL stores saved searches, job metadata, and
chat sessions. Source CSVs, staging files, and index snapshots live in object
storage. The OpenAPI schema is rendered as a static reference on GitHub Pages.

Indices are versioned. A new index is built as \texttt{papers\_vN}, smoke-tested
against a fixed query set, and then made visible through an alias swap. Cache keys
include the alias version, so a reindex cannot serve stale results. Kafka is
deliberately deferred: at this corpus size a CLI job is simpler, and the worker
entry point already exists for the day replay or multiple consumers justify a
broker.

\section{Search and ranking}

The index stores one document per measurement property. Each document carries a
single property name, value, and unit, which makes terms aggregations and range
filters straightforward. Documents are grouped by \texttt{proxy\_id} for display
and by DOI for paper pages.

Text retrieval uses a multi-field BM25 query with boosts on the composition
label, property name, application, and journal, plus a catch-all field, following
the retrieval and aggregation model described by \citet{gormley2015}. Elements
are indexed as keywords, so ``Al and Cr'' is a filter rather than a text match.
Numeric filters apply to \texttt{property.value}, and the active property
determines the range slider. Filtering follows OR within a field and AND across
fields, matching the expectations of faceted search interfaces.

Pagination uses \texttt{search\_after} with an opaque cursor. Deep pagination is
not cached; only the first page is. Aggregations for charts are cached
separately and refreshed on ingest. Phase two adds composition similarity: a
118-dimensional, periodic-table-ordered vector supports a k-nearest-neighbour
query, which is chemically meaningful and avoids text embeddings entirely.

\section{Per-paper chat}

The chat interface answers questions about one paper at a time. Its context is
the paper's measurement rows rather than extracted PDF text. For typical papers
the rows fit in a single prompt; for larger ones the model receives read-only
tools to query, aggregate, and fetch records. Every numeric answer must cite a
\texttt{proxy\_id} and table number, and the prompt instructs the model to say
that a value is absent rather than infer one.

This design avoids a chunk index, an embedding pipeline, and a vector database.
It also sets a clear expectation: the assistant can compare and summarise the
dataset, but it cannot explain a measurement method or a heat treatment that the
corpus does not record.

\section{Evaluation plan}

Evaluation uses three complementary methods. First, a known-item set of queries
built from held-out records measures recall and mean reciprocal rank. Second, a
graded relevance set for faceted queries measures nDCG at ten. Third, operational
metrics cover p95 latency, facet computation time, and cache hit rate. The
ingestion pipeline is evaluated separately through rejection rate and reason
distribution, because silent cleaning failures are more damaging than slow
queries.

Chat quality is assessed on a small question set per paper with three outcomes:
correct and cited, correct but uncited, and unsupported. Unsupported answers are
the primary failure mode and are tracked explicitly.

\section{Conclusion}

OMAT-SE is a search engine for measurement records, not a document search engine
with a materials theme. Treating composition and property as first-class,
filterable entities produces a simpler index, cheaper facets, and a chat
interface that can be grounded without a retrieval pipeline. The remaining work
is implementation: the canonical models, the cleaning pipeline, the index
mapping, and the search handlers.

\bibliographystyle{plainnat}
\bibliography{references}

\end{document}
